\documentclass[../../../include/open-logic-section]{subfiles}

\begin{document}

\olfileid{sth}{choice}{hartogs}
\olsection{তুলনীয়তা ও হার্টগসের লেমা}

এটি সুবিধার দিক। অসুবিধাও আছে: চয়ন না ধরলে
বিষয়টি \emph{জটিল} হয়ে পড়ে। কেন, তা বোঝাতে
\cite{Hartogs1915}-এর একটি ফল দেখি:

\begin{lem}[\emph{$\ZF$-এ}]\ollabel{HartogsLemma}
যেকোনো সেট $A$-র জন্য এমন একটি ক্রমসংখ্যা
$\alpha$ আছে যে $\cardnless{\alpha}{A}$।
\end{lem}

\begin{proof}
$B \subseteq A$ এবং $R \subseteq B^2$ হলে
\olref[ord-arithmetic][using-addition]{rankcomputation}
অনুযায়ী $\tuple{B, R} \subseteq
V_{\setrank{A}+4}$। তাই পৃথকীকরণ দিয়ে নিই:
\[
	C = \Setabs{\tuple{B, R} \in V_{\setrank{A}+5}}{B\subseteq A
	\text{ and $\tuple{B, R}$ is a well-ordering}}
\]
প্রতিস্থাপন এবং
\olref[ordinals][ordtype]{thmOrdinalRepresentation}
ব্যবহার করে সেটটি গড়ি:
\[
	\alpha = \Setabs{\ordtype{B, R}}{\tuple{B, R} \in C}.
\]
\olref[ordinals][basic]{corordtransitiveord}
অনুযায়ী $\alpha$ একটি ক্রমসংখ্যা, কারণ এটি
ক্রমসংখ্যার পরিযায়ী সেট। সত্যিই,
$\gamma \in \beta \in \alpha$ হলে কোনো
% BN-SRC-463: B is a subset of A, not of the relation R.
$B \subseteq A$-র জন্য
$\beta = \ordtype{B, R}$। তখন
\olref[ordinals][iso]{wellordinitialsegment}
অনুযায়ী কোনো $b\in B$-র জন্য
$\gamma = \ordtype{B_b, R_b}$; সুতরাং
$\gamma \in \alpha$।

বিরোধের জন্য ধরি,
$f \colon \alpha \to A$ একটি !!a{injection}।
তাহলে নিই:
% BN-SRC-464: bind two elements of the ordinal domain, not the ordinal itself.
\begin{align*}
	B &= \ran{f}\\
	R &= \Setabs{\tuple{f(\xi), f(\eta)} \in A \times A}{\xi,\eta\in\alpha\text{ and }\xi\in\eta}.
\end{align*}
স্পষ্টত $\alpha = \ordtype{B, R}$ এবং
$\tuple{B, R} \in C$। ফলে $\alpha\in\alpha$,
যা বিরোধ।
\end{proof}

এ থেকে একটি গভীর ফল পাওয়া যায়:

\begin{thm}[\emph{$\ZF$-এ}]
নিচের দাবিদুটি সমতুল্য:
\begin{enumerate}
	\item\ollabel{equivwo} সুক্রমবিন্যাসের স্বতঃসিদ্ধ
	\item\ollabel{equivcompare} যেকোনো সেট $A$ ও
	$B$-র জন্য হয় $\cardle{A}{B}$, নয়
	$\cardle{B}{A}$
\end{enumerate}
\end{thm}

\begin{proof}
\emph{\olref{equivwo} $\Rightarrow$
\olref{equivcompare}।} $A$ ও $B$ স্থির করি।
\olref{equivwo} অনুযায়ী $\tuple{A, R}$ ও
$\tuple{B, S}$-র সুক্রমবিন্যাস আছে।
\olref[ordinals][ordtype]{thmOrdinalRepresentation}
থেকে ক্রমসমরূপতা নিই
% BN-SRC-465: isomorphisms map ordinals to the carrier sets A and B.
$f \colon \alpha \to A$ এবং
$g \colon \beta \to B$।
\olref[sth][ordinals][basic]{ordinalsaresubsets}
অনুযায়ী হয় $\alpha \subseteq \beta$,
নয় $\beta \subseteq \alpha$।
$\alpha\subseteq\beta$ হলে
$\comp{f^{-1}}{g} \colon A \to B$
একটি !!a{injection}, তাই
$\cardle{A}{B}$; $\beta\subseteq\alpha$ হলে একইভাবে
$\cardle{B}{A}$।

\emph{\olref{equivcompare} $\Rightarrow$
\olref{equivwo}।} $A$ স্থির করি।
\olref{HartogsLemma} অনুযায়ী এমন ক্রমসংখ্যা
$\beta$ আছে যে $\cardnless{\beta}{A}$।
\olref{equivcompare} থেকে
$\cardle{A}{\beta}$। ফলে একটি !!{injection}
$f \colon A \to \beta$ আছে। এর সাহায্যে
$A$-র সদস্যগুলিকে সুক্রমবিন্যস্ত করতে
$\Setabs{\tuple{a, b} \in A \times A}{f(a)
\in f(b)}$ সম্পর্কটি নিই।
\end{proof}
\noindent
অতএব সুক্রমবিন্যাস ব্যর্থ হলে কিছু সেটের
আকার \emph{আক্ষরিক অর্থেই অতুলনীয়}। তখন
ট্রান্সফাইনাইট অঙ্কবাচক পাটিগণিতও জটিল হবে।
যেমন, অসীম $A$ ও $B$-র জন্য বিযুক্ত যোগফল
ও কার্তেসীয় গুণফল উভয়ই তাদের মধ্যে বৃহত্তর
$M$-এর সমসংখ্যক—আগের এই ধারণা
(দেখো \olref[card-arithmetic][simp]{cardplustimesmax})
সার্বিকভাবে বলা যাবে না।
% BN-SRC-466: the source's nested cardeq expression is ill-typed;
% spell out the two intended size comparisons separately.
$\cardeq{A \disjointsum B}{M}$ এবং
$\cardeq{A \times B}{M}$ বলতেই আগে
$A$ ও $B$-র আকার তুলনা করে কোনটি বৃহত্তর,
তা স্থির করা চাই। তুলনাই অসম্ভব হলে
``কোনটি বড়'' প্রশ্নটির অর্থ নেই।

\end{document}
